\documentclass[11pt, letterpaper]{article}
\usepackage[utf8]{inputenc}
\usepackage[margin=1.5cm]{geometry}
\usepackage{array}
\usepackage{booktabs}
\usepackage{enumitem}
\usepackage{hyperref}
\hypersetup{
  colorlinks=true,
  linkcolor=blue,
  urlcolor=blue
}

\title{\textbf{A Reproducible Pipeline for Classifying Bug Types in LLM-Generated Code}}
\author{Zixiao Xie\\
School of Information Sciences, University of Illinois Urbana-Champaign\\
\texttt{zixiaox2@illinois.edu}}
\date{}

\begin{document}
\maketitle

\noindent
\textbf{Abstract}\\
This software provides a reproducible pipeline for studying bug types
in code generated by large language models (LLMs). Given a benchmark such as
HumanEval and a set of model generations, the pipeline extracts runnable code,
executes tests in a subprocess harness, captures failure signals, assigns
bug-category labels, refines wrong-output LOGIC errors into subtypes, computes automatic
agreement statistics, and produces tables, figures, and manuscript-ready
artifacts. The current release analyzes 656 generations from four hosted
models and supports reproducible empirical studies that go beyond aggregate pass@k
scores.

\vskip0.4cm
\noindent
\textbf{Keywords}\\
large language models; code generation; software bugs; empirical software
engineering; reproducibility; research software

\newpage
\noindent
\textbf{Code metadata}\\

\vskip0.2cm
\noindent
\begin{tabular}{|p{0.8cm}|p{6.2cm}|p{9.8cm}|}
\hline
\textbf{Nr.} & \textbf{Code metadata description} & \textbf{Project metadata} \\
\hline
C1 & Current code version & 0.1.0 \\
\hline
C2 & Permanent link to code/repository used for this code version & TBD: public repository URL \\
\hline
C3 & Permanent link to Reproducible Capsule & TBD: archived release DOI or reproducible capsule URL \\
\hline
C4 & Legal Code License & MIT License \\
\hline
C5 & Code versioning system used & Git \\
\hline
C6 & Software code languages, tools, and services used & Python, LaTeX, OpenRouter API, OpenAI-compatible chat-completions API \\
\hline
C7 & Compilation requirements, operating environments \& dependencies & Python 3.10+; pandas, matplotlib, scipy, pyyaml, openai; optional LaTeX installation for manuscript builds. Tested on Windows; expected to run on Linux/macOS with the same Python dependencies. \\
\hline
C8 & If available Link to developer documentation/manual & README.md, PROTOCOL.md, and submission\_materials/ in the repository \\
\hline
C9 & Support email for questions & \texttt{zixiaox2@illinois.edu} \\
\hline
\end{tabular}

\section*{Software Purpose and Functionality}
LLM code-generation evaluation is often summarized by aggregate functional
correctness metrics such as pass@k. Those metrics leave a
question unanswered: when generated code fails, what kind of bug is it?
This software addresses that question by turning raw model generations into an
inspectable failure-analysis dataset. It is designed for empirical software
engineering researchers, LLM evaluation researchers, educators, and practitioners
who want to inspect not only whether generated code passes tests, but how it fails.

The pipeline consists of eight stages. First, it loads HumanEval problems and
generates one or more model responses through an OpenAI-compatible API. Second, it
extracts runnable Python functions from chat-style responses. Third, it executes
each generation against the benchmark tests in a subprocess harness with timeout
protection. Fourth, it records pass/fail status and failure signals. Fifth, it
labels failures into primary bug categories using deterministic rules for syntax,
timeout, and exception failures, and LLM-assisted classification for wrong-output
assertion failures. Sixth, it refines the largest observed LOGIC category into seven
subtypes. Seventh, it runs an independent blind LLM annotator and computes
LLM-vs-LLM agreement. Eighth, it prepares a deterministic blind human-review sheet
for optional manual validation. Finally, it generates tables, figures, and
arXiv-ready manuscript artifacts.

\section*{Illustrative Results}
The current release applies the pipeline to four instruction
models: GPT-4o-mini, DeepSeek-Chat, Qwen2.5-72B-Instruct, and
Llama-3.1-8B-Instruct. Model identifiers, provider endpoint, access date,
generation settings, and labeler settings are recorded in the artifact file
\texttt{results/tables/model\_metadata.csv}. Across 656 HumanEval generations, 546 pass and 110 fail.
Among the failures, 87.3\% are wrong-output LOGIC errors: the code executes but returns a
wrong result (95\% Wilson CI: 79.8--92.3\%). A second-stage subtype analysis of 96
LOGIC failures finds that mathematical-reasoning errors (44.8\%, 43/96; 95\% CI:
35.2--54.7\%) and wrong-algorithm choices (21.9\%, 21/96; 95\% CI: 14.8--31.1\%)
are the largest categories. A category-by-model chi-square test does not detect an
association ($\chi^2=21.290$, $p=0.265$, Cramer's $V=0.254$), but the contingency table is
sparse, so the comparison is reported as exploratory; a deterministic 10{,}000-iteration
permutation check gives $p=0.2691$. A lower-dimensional
LOGIC/non-LOGIC companion test also does not detect an association ($\chi^2=3.394$,
$p=0.335$, Cramer's $V=0.176$; permutation $p=0.3445$).

The release also includes a robustness check for the fine-grained subtype labels.
An independent blind LLM annotator labels all 96 LOGIC failures without seeing the
reference solution, the primary subtype label, or the primary explanation. Exact
agreement is 42/96 (43.8\%), with Cohen's $\kappa=0.305$, conventionally
interpreted as fair agreement. This result is reported as a limitation: the
software records label uncertainty rather than omitting it.

\section*{Impact Overview}
The software packages this LLM code-failure analysis as a reproducible and inspectable
artifact.
Researchers can rerun the full analysis, audit individual failure
labels, inspect disagreement cases, regenerate figures, and adapt the taxonomy for
new studies. The software also connects benchmark-style evaluation with
qualitative bug analysis: it keeps execution-based unit-test results while adding
interpretable failure categories.

The pipeline can be used for several research questions. It can be used to study
whether higher-pass-rate models fail less or fail differently, whether
reasoning-related bugs account for more failures than surface syntax/API mistakes, how stable
LLM-assisted bug labels are under independent annotation, and which failure types
should be prioritized by repair or verification tools. Because the intermediate
artifacts are stored as JSONL/CSV files, the software can be extended without
re-running every stage.

\section*{Reuse Potential}
The current implementation targets Python and HumanEval, but the design is
modular. Researchers can replace the benchmark loader, add MBPP or repository-level
tasks, compare additional models, increase the number of samples per problem,
replace OpenRouter with another OpenAI-compatible endpoint, or substitute human
annotation for the automatic labeling stages. The agreement-analysis utilities can
also be reused for studies that compare human labels, LLM labels, or hybrid
labeling workflows. The current release includes a 40-item blind human-review
sample with blank human-label fields; it is provided as a review instrument, not
as evidence of completed human validation.

Educators can use the generated tables and disagreement files to teach common
failure modes in AI-generated code. Tool builders can use the failure categories
to prioritize semantic testing, edge-case generation, and repair strategies. The
pipeline is intentionally small and script-based, with intermediate files stored
as JSONL and CSV rather than hidden inside a benchmark service.

\section*{Limitations and Future Improvements}
The current empirical results are limited to HumanEval, Python, one sample per
problem, and hosted model identifiers that may route to provider-current weights.
Fine-grained subtype labels are machine-assigned, and the independent agreement
check shows only fair agreement. A blind human-review sheet is included, but
human labels have not yet been completed. Future releases should add completed
human validation, more samples per task, additional benchmarks such as MBPP,
non-Python languages, and repairability experiments by bug type. Public repository
and archived-release links must be inserted before formal submission.

\section*{Acknowledgements}
No acknowledgements.

\section*{References}
\begin{enumerate}[leftmargin=*]
  \item Z. Xie, ``A Reproducible Pipeline for Classifying Bug Types in LLM-Generated Code,'' version 0.1.0, archived software release, TBD: archived software DOI.
  \item M. Chen et al., ``Evaluating Large Language Models Trained on Code,'' arXiv:2107.03374, 2021.
  \item J. Austin et al., ``Program Synthesis with Large Language Models,'' arXiv:2108.07732, 2021.
  \item J. Liu, C. S. Xia, Y. Wang, and L. Zhang, ``Is Your Code Generated by ChatGPT Really Correct? Rigorous Evaluation of Large Language Models for Code Generation,'' NeurIPS, 2023.
  \item F. Tambon, A. Nikanjam, F. Khomh, and G. Antoniol, ``Bugs in Large Language Models Generated Code: An Empirical Study,'' Empirical Software Engineering, 2025.
  \item J. Cohen, ``A coefficient of agreement for nominal scales,'' Educational and Psychological Measurement, 1960.
\end{enumerate}

\end{document}
